\documentclass[oneside,11pt,openright]{report}

\usepackage[latin1]{inputenc}
\usepackage[american]{babel}
\usepackage{a4}
\usepackage{latexsym}
\usepackage{amssymb}
\usepackage{amsmath}
\usepackage{epsfig}
\usepackage[T1]{fontenc}
\usepackage{lmodern}
\usepackage[labeled]{multibib}
\usepackage{color}
\usepackage{datetime}
\usepackage{epstopdf} 
\usepackage{graphicx}
\usepackage{pgfplots}
\usepackage[format=hang]{caption}

\usepackage{pgf,tikz}
\usepackage{comment}
\usetikzlibrary{arrows,automata}
\usetikzlibrary{backgrounds,fit}
\usetikzlibrary{shapes,patterns}
\usetikzlibrary{calc,chains,positioning}

\renewcommand*\ttdefault{txtt}
\newcommand{\BigO}[1]{\ensuremath{\operatorname{O}\left(#1\right)}}
\newcommand{\BigT}[1]{\ensuremath{\Theta\left(#1\right)}}
\newcommand{\specialcell}[2][c]{%
  \begin{tabular}[#1]{@{}c@{}}#2\end{tabular}}
% see http://imf.au.dk/system/latex/bog/

\newcommand{\adjustimg}{% Horizontal adjustment of image
  \ifodd\value{page}\hspace*{\dimexpr\evensidemargin-\oddsidemargin}\else\hspace*{-\dimexpr\evensidemargin-\oddsidemargin}\fi%
}
\newcommand{\centerimg}[2][width=\textwidth]{% Center an image
  \makebox[\textwidth]{\adjustimg\includegraphics[#1]{#2}}%
}
\newcommand{\MakeHeap}{\textsc{MakeHeap}}
\newcommand{\FindMin}{\textsc{FindMin}}
\newcommand{\Insert}{\textsc{Insert}}
\newcommand{\DeleteMin}{\textsc{DeleteMin}}
\newcommand{\DecreaseKey}{\textsc{DecreaseKey}}
\newcommand{\Delete}{\textsc{Delete}}
\newcommand{\Meld}{\textsc{Meld}}
\newcommand{\Dijkstra}{\textsc{Dijkstra}}
\newcommand{\NULL}{\textbf{null}}

\newcommand{\Member}{\textsc{Member}}
\newcommand{\Minimum}{\textsc{Minimum}}
\newcommand{\Maximum}{\textsc{Maximum}}
\newcommand{\Predecessor}{\textsc{Predecessor}}
\newcommand{\Successor}{\textsc{Successor}}

\newcommand{\HIGH}{\textsf{high}}
\newcommand{\LOW}{\textsf{low}}
\newcommand{\INDEX}{\textsf{index}}
\newcommand{\HIGHER}{\sqrt[\uparrow]{u}}
\newcommand{\LOWER}{\sqrt[\downarrow]{u}}

\begin{document}

%%%%%%%%%%%%%%%%%%%%%%%%%%%%%%%%%%%%%%%%%%%%%%%%%%%%%%%%%%%%%%%%%%%%%%%

\pagestyle{empty} 
\pagenumbering{roman} 
\vspace*{\fill}\noindent{\rule{\linewidth}{1mm}\\[4ex]
{\Huge\sf Binary Heaps, Fibonacci Heaps and\\[2ex]Dijkstras shortest path}\\[4ex]
{\huge\sf Kristoffer Just Andersen, 20051234\\[2ex]
\huge\sf Troels Leth Jensen, 20051234 \\[2ex]
\huge\sf Morten Krogh-Jespersen, 20022362}\\[2ex]
\noindent\rule{\linewidth}{1mm}\\[4ex]
\noindent{\Large\sf Project 1, Advanced Data Structures 2013, Computer Science\\[1ex] 
\monthname\ \the\year  \\[1ex] Advisor: Gerth Stølting Brodal\\[15ex]}\\[\fill]}
\epsfig{file=logo.eps}\clearpage

%%%%%%%%%%%%%%%%%%%%%%%%%%%%%%%%%%%%%%%%%%%%%%%%%%%%%%%%%%%%%%%%%%%%%%%

\tableofcontents
\pagenumbering{arabic}
\setcounter{secnumdepth}{2}

%%%%%%%%%%%%%%%%%%%%%%%%%%%%%%%%%%%%%%%%%%%%%%%%%%%%%%%%%%%%%%%%%%%%%%%

\chapter{Introduction}

needs content

\chapter{Bit-vector}

A  bit-vector  or  a  bit  array  is  a  datastructure  that  supports
$\Insert$,      $\Delete$      and     $\Member$      in      constant
time~\cite[p. 532]{ITA09}. The  bit-vector consists of a  set of bits,
for  which  the  $i'th$  bit  indicates the  presence  of  the  $i'th$
key. Thus, in order to support queries over an universe of size $u$ we
must have a bit-vector with $u$ bits.

To $\Insert$ some key $i$ into the bit-vector the $i'th$ bit is set to
true,  similarly to  $\Delete$ the  same key  $i$, the  $i'th$ bit  is
cleared. The $\Member$ test for key $i$ is done by checking the $i'th$
bit.   All  these   operations  are   supported  easily   through  bit
shifts~\footnote{\begin{align*}   &    \textsc{TEST\_BIT}(var,pos)   =
    ((var)\quad \&=\quad(1<<(pos))) \\  & \textsc{SET\_BIT}(var,pos) =
    ((var)\quad|=\quad(1<<(pos))) \\  & \textsc{CLEAR\_BIT}(var,pos) =
    ((var)\quad\&=\quad\textsf{\textasciitilde}(1<<(pos)))
\end{align*}}. 
Please note, that  these bit operations assume that  the universe size
does  not  exceed the  word  size.  If that  is  not  the case,  extra
computation is needed to first find the correct block; this however is
also constant time.

The   bit-vector   performs    badly   for   $\Minimum$,   $\Maximum$,
$\Predecessor$  and $\Successor$  because  the entire  array could  be
traversed in search for the answer.

\begin{center}
  \begin{tabular}{ l | c }
    Operation & Bit-Vector  \\ \hline
    \Member & $\BigT{1}$ \\ 
    \Minimum & $\BigT{u}$ \\ 
    \Maximum & $\BigT{u}$ \\ 
    \Predecessor & $\BigT{u}$ \\ 
    \Successor & $\BigT{u}$ \\ 
    \Insert & $\BigT{1}$ \\ 
    \Delete & $\BigT{1}$ \\
  \end{tabular}
\end{center}

Alternatively, $\Minimum$ and $\Maximum$  can be maintained during all
changes to  the bit vector, i.e.  during insert and delete.  With this
approach the running  time $\Delete$ would in worst  case be completly
dominated  by the  linear search  for a  new $\Minimum$  or $\Maximum$
while querying for the $\Minimum$  or $\Maximum$ would become constant
time operations.

\chapter{Van Emde Boas Tree}

A  van Emde  Boas tree  is a  recursive data  structure that  supports
finding the predecessor and successor in \BigT{\log \log m} time where
$m$  is the  size  of the  universe,  in other  words,  the amount  of
distinct    keys   with    a    total   ordering    that   the    tree
support~\cite[p. 545]{ITA09}~\cite{VEB04}. Because van Emde Boas Trees
allow  universe  sizes  of  any  power of  2,  we  denote  $\HIGHER  =
2^{\lceil{(\log   u)/2}\rceil}$   and   $\LOWER   =   2^{\lfloor{(\log
    u)/2}\rfloor}$.   Each level  of  the  van Emde  Boas  Tree has  a
universe $u$ and  it contains $\HIGHER$ clusters/bottom  van Emde Boas
trees of universe size $\LOWER$ and one auxilary van Emde Boas Tree of
size $\HIGHER$ we denote top.  For each recursion the universe shrinks
by $\sqrt{u}$.

For each  tree we have two  attributes that store the  minimum and the
maximum key.  The minimum  key cannot  be found in  any of  the bottom
trees, neither  can the  maximum unless it  differs from  the minimum,
which happens  if there  are more  than one element  in the  tree. The
attributes helps reduce the number of recursive calls, because one can
in constant time decide if a  value lies within the range, without the
need to recurse.

Van Emde Boas Trees utilizes that  keys are represented as bits and if
we view a key $x$ as a  $\log u$-bit binary integer, we can divide the
bits up in  a most significant and a least  significant part. The most
significant   part  of   the   bit  vector   identifies  the   cluster
$\lfloor{x/\sqrt{u}}\rfloor$ where $x$ will  apear in position $x \mod
\sqrt{u}$ in the beforementioned bottom  tree. Let us denote $\HIGH(x)
= \lfloor{x/\sqrt{u}}\rfloor$ and $\LOW(x) =  x \mod \sqrt{u}$ then we
get the identity $x = \HIGH(x) \sqrt{u} + \LOW(x)$.

Below  is  the time-complexities  of  each  operation listed  for  the
corresponding  data structure  and  let  $u$ denote  the  size of  the
universe and $n$ the number of keys stored in the structure:

\begin{center}
  \begin{tabular}{ l | c | c | c }
    Operation & Bit-Vector & Red-Black Tree & van Emde Boas Tree \\ \hline
    \Member & $\BigT{1}$ & $\BigT{\log n}$ & $\BigO{\log \log u}$ \\ 
    \Minimum & $\BigT{u}$ & $\BigT{\log n}$ & $\BigO{1}$\\ 
    \Maximum & $\BigT{u}$ & $\BigT{\log n}$ & $\BigO{1}$ \\ 
    \Predecessor & $\BigT{u}$ & $\BigT{\log n}$ & $\BigO{\log \log u}$  \\ 
    \Successor & $\BigT{u}$ & $\BigT{\log n}$ & $\BigO{\log \log u}$ \\ 
    \Insert & $\BigT{1}$ & $\BigT{\log n}$ & $\BigO{\log \log u}$ \\ 
    \Delete & $\BigT{1}$ & $\BigT{\log n}$ & $\BigO{\log \log u}$ \\
  \end{tabular}
\end{center}

\section{Finding minimum or maximum key}

This is  a constant operation, since  the van Emde Boas  Tree directly
stores the minimum and maximum element.

\section{Finding a member}

Finding  out  if  a  key  is  stored  in  a  van  Emde  Boas  Tree  is
simple. Either the  key is the minimum or the  maximum element or else
recurse until we find the element. It is easy to figure out which tree
to recurse on, since this is  the $\HIGH(x)$-th tree, and we just have
to search in the smaller tree for $\LOW(x)$.

So how long does the search take in worst-case? The data structure can
be described by the following recurrence function:
\begin{align*}
    T(u) \leq T(\HIGHER) + \BigO{1}
\end{align*}
If we let $m = \log  u$ and realize that $\lceil{m/2}\rceil \leq 2m/3$
for $m \geq 2$,  which is the leaf size of the van  Emde Boas Tree, we
get:

\begin{align*}
    T(u)  & \leq T(\HIGHER) + \BigO{1} \\
    \Downarrow \\
    T(2^m) & \leq T(2^{\lceil{m/2}\rceil}) + \BigO{1} \\
    \Downarrow \\
    T(2^m) & \leq T(2^{2m/3}) + \BigO{1} \\
    \Downarrow \\
    S(m) & \leq S(2m/3) + \BigO{1} \\
\end{align*}

By the  master thorem~\cite[p.  93]{ITA09}, has  the solution  $S(m) =
\BigO{\log  m}$. Because  $T(u)  =  T(2^m)$ we  get  $\BigO{\log m}  =
\BigO{\log \log  u}$. Therefore  the procedure takes  $\BigO{\log \log
  u}$.

\section{Finding a successor or a predecessor}

As with finding a member, for the base case a successor or predecessor
can be found in  constant time. If not, we have  to determine where to
find  the next  element we  are searching  for. Finding  successor and
predecessor is  completely analog  thus we  only describe  finding the
successor for an element $x$.

First,  we check  if  the sucessor  is  in the  bottom  tree at  index
$\HIGH(x)$. This  check ban  be done  in constant time  since it  is a
matter of checking  maximum for that tree. If maximum  exist and it is
higher  than  $\LOW(x)$,  we  know  we  have  to  search  inside  that
particular tree. This will take $\BigO{\log\log u}$ time.

If maximum does not exist or it  is less than $\LOW(x)$ we use the top
tree to search for the successor to $\HIGH(x)$. If we can find such an
element it will gives  us an index to a bottom  tree $b$. Hereafter it
is a  constant lookup to  find $\Minimum(b)$.  This case also  runs in
$\BigO{\log\log u}$ time since we have on search in top and a constant
lookup. Therefore, the total running time of successor and predecessor
is $\BigO{\log\log u}$.

\section{Insert}

Insert is pretty simple. One of the following can happen:

\begin{itemize}
\item The list  is empty, which can be discovered  in constant time by
  checking the minimum attribute. If that is the case, insert just set
  min and max.
\item If not in the base case, find out if the bottom of $\HIGH(x)$ is
  empty. If that is the case, we can just set min and max as above and
  update the top structure. If not,  we just call insert on the bottom
  tree for $\LOW(x)$.
\end{itemize}

Always remember to  set max or swap  with min (because min  was not in
the tree, but now has to be  since it would no longer be min). Either,
the insert function  recurses on the top tree of  size $\HIGHER$ or on
one of the  bottom trees of size $\LOWER$ but  not on both. Therefore,
the running time is at most $\BigO{\log \log u}$.

\section{Delete}

If there is only one element or  universe size is 2 is easy to perform
delete  in  constant  time.  Otherwise,  some  more  work  has  to  be
done. Consider  what should be  done if we  try to delete  the minimum
element. Since  min is  not stored in  the tree, we  have to  find the
minimum element in the bottom trees and make sure we delete it so that
it can be placed as the new min. Of course, it can happen that the the
bottom tree becomes  empty so the top structure has  to be updated. It
might  also be  that we  remove  max but  updating max  is a  constant
operation.

In the description  above we actually could make  two recursive calls;
one to update the bottom tree and  possibly one to update the top. But
if  we update  the top  then the  bottom tree  would only  contain one
element. But  if that is the  case the first recursive  function takes
constant time. Therefore delete runs in $\BigO{\log \log n}$.

\chapter{Combining vEB tree and bit vectors}

As mentioned in TODO, when $u$ becomes suffitiantly small it can be
benifial to switch to a bit-vector or red black tree. This is possible
because of the recursive nature of vEB trees. To decide membership of
elements over some universe it is done by deciding it recursively over
smaller universes until some minimum size is reached, thus it is no
problem to switch to a different datastructure that enables the same
queries as vEB trees for universes small enough.

We have implemented as vEB tree that allows switching to bit-vectors
when the universe falls below some constant, in our testing we will
examine what the value of this constant means for the performance of
our implementation.

With some constant in mind, we can conside the total universe
$U_{total} (0....2^{24}-1 for this handin)$ as the product of two
different universes $U_{vEB}$ and $U_{bit}$, respectivly the universes
decided exclusively by the vEB tree and bit vectors. Now the vEB tree
that uses bit vectors for leafs doesn't strictly speak decide
membership of the elements, but they decide which bit-vector would
contain the element if it is present.

\begin{center}
  \begin{tabular}{ l | c }
    Operation & Combi-tree  \\ \hline
    \Member & $\BigO{ \log \log u_{veb} }$ \\ 
    \Minimum & $\BigO{1}$ \\ 
    \Maximum & $\BigO{1}$ \\ 
    \Predecessor & $\BigO{ u_{bit} + \log \log u_{veb} }$ \\ 
    \Successor & $\BigO{ u_{bit} + \log \log u_{veb} }$ \\ 
    \Insert & $\BigO{\log \log u_{veb} }$ \\ 
    \Delete & $\BigO{\log \log u_{veb} }$ \\
  \end{tabular}
\end{center}



\chapter{Conlusion}

needs contents


%%%%%%%%%%%%%%%%%%%%%%%%%%%%%%%%%%%%%%%%%%%%%%%%%%%%%%%%%%%%%%%%%%%%%%%

\addcontentsline{toc}{chapter}{Bibliography}
\bibliographystyle{plain} 
\bibliography{report}

\end{document}

